
\chapter*{前言}
\addcontentsline{toc}{chapter}{前言}
\markboth{前言}{前言}

随着米哈伊尔·格罗莫夫（Mikhail Gromov）在 20 世纪 80 年代发展出伪全纯
曲线方法，以及安德烈亚斯·弗洛尔（Andreas Floer）随后创立了导向弗洛尔
上同调的椭圆型莫尔斯理论，辛几何的面貌发生了巨大变化。辛几何中许多从前
难以处理的问题，都由伪全纯曲线技术解决；辛拓扑的概念也逐渐成形。与此同
时，物理学中也出现了平行的发展，其中尤以闭弦理论为著。

1993 年，第一署名作者深谷贤治（Kenji Fukaya；日文 深谷賢治）部分受到
西蒙·柯万·唐纳森（Simon Kirwan Donaldson）在弗洛尔上同调中构造裤形
乘积的启发，把 $A_\infty$ 范畴结构引入辛几何；其中的对象是拉格朗日子
流形，态射则是弗洛尔上同调（或复形）。以这一代数框架为基础，
马克西姆·孔采维奇（Maxim Kontsevich）在 1994 年苏黎世国际数学家大会的演讲中，
提出了著名的同调镜像对称：凝聚层导出范畴与拉格朗日子流形的深谷范畴之间
的对称。随着后来开弦理论中 D-膜研究的发展，这一同调镜像对称得到进一步
丰富，并成为代数几何、辛几何以及物理学中许多新见解与新进展的源泉。不过，
同调镜像对称尚未得到严格表述，主要原因是人们对于拉格朗日子流形本身的
弗洛尔理论仍缺乏了解。

本书说明构造弗洛尔上同调时障碍与反常如何产生，给出这些障碍的精确定式，
进而对弗洛尔上同调的形变理论作细致的代数与分析研究。事实表明，即使只是
以数学上精确的方式描述这类障碍，也需要一套关于滤过 $A_\infty$ 代数的新
同调代数。此外，要在拉格朗日子流形的几何语境中验证这种代数结构的存在性，
还须对相应的伪全纯圆盘模空间作非平凡的分析研究。我们也把如此构造的弗洛尔
上同调立即应用于辛拓扑中的各种问题。其中许多结果改进了此前由弗洛尔理论
得到的已知结果，若干对同调镜像对称的直接应用也是新的。我们期待，随着这套
理论的真正潜力逐步显露并付诸实现，更多非平凡的应用很快会随之出现。

在我们撰写本书期间，辛几何及相关领域中出现了若干重要进展。拓扑弦、D-膜、
伪全纯曲线与辛几何中的弗洛尔理论之间的关系，如今已经理解得更加清楚。我们
在写作本书的过程中一直倡导在弗洛尔理论中使用高阶代数结构；如今，这已经
成为一个热门而不可或缺的研究方向。不仅如此，处理非线性偏微分方程解空间之
各类模空间的技术也有所进步，并同高阶代数结构的形式体系交织在一起，使几何
图景变得更加清晰。这将有助于几何理论的进一步发展。本书充分利用这些进展，
以目前所能达到的最一般形式给出拉格朗日子流形的弗洛尔理论。我们希望，本书
能成为辛几何与同调镜像对称未来进展的一块踏脚石。

最终完成本书的这次合作，始于 1996 年 7 月 8 日至 12 日在瑞士阿斯科纳举行
的一次会议。当时我们几乎没有想到，这一项目竟会持续十余年。我们从这次合作
中获得了莫大的乐趣，也希望在未来几十年继续合作。事实上，我们在辛拓扑与
镜像对称这片焕然一新的领域中已经开始第二段旅程，并期待收获更多合作成果：
眼前的景象与十三年前留在身后的景象迥然不同，也振奋人心得多！

\begin{flushright}
2009 年 6 月 9 日\par
深谷贤治（Kenji Fukaya）\quad 吴龙根（Yong-Geun Oh）\par
太田啓史（Hiroshi Ohta）\quad 小野薫（Kaoru Ono）
\end{flushright}
